\editionPage[\pdfbookmark[0]{Corrigés des défis}{bm-defi-cor-1-parent}]{defi-cor-1}{D01 · Corrigé}{Une fraction entre deux autres}{\hyperref[defi-ex-1]{Énoncé p.\,\pageref*{defi-ex-1}}}
\begin{blocEdition}{defi-cor-1}
$5/12=25/60$ et $7/15=28/60$. Le numérateur doit être strictement entre 25 et 28 : on obtient $\res{26/60\text{ et }27/60}$, soit exactement deux fractions demandées. Au dénominateur 120, les bornes sont 50/120 et 56/120 ; par exemple $\res{53/120}$ convient. Il existe cinq choix de numérateur, de 51 à 55. Attention : compter des écritures de dénominateur imposé ne signifie pas qu’elles soient irréductibles.
\end{blocEdition}
\editionPage{defi-cor-2}{D02 · Corrigé}{Comparer sans grands calculs}{\hyperref[defi-ex-2]{Énoncé p.\,\pageref*{defi-ex-2}}}
\begin{blocEdition}{defi-cor-2}
$2^{10}=1024>1000=10^3$. En élevant au cube des nombres positifs, $2^{30}=1024^3>1000^3=10^9$. En élevant cette dernière inégalité au carré, $\res{2^{60}>10^{18}}$. Une valeur approchée n’est pas nécessaire. La comparaison utilise seulement les règles des puissances et la croissance du carré et du cube sur les nombres positifs.
\end{blocEdition}
\editionPage{defi-cor-3}{D03 · Corrigé}{Une recherche entière exhaustive}{\hyperref[defi-ex-3]{Énoncé p.\,\pageref*{defi-ex-3}}}
\begin{blocEdition}{defi-cor-3}
$180=2^2\times3^2\times5$. Un diviseur est de la forme $2^a3^b5^c$, avec $0\le a,b\le2$ et $0\le c\le1$. Pour être un carré, les exposants doivent être pairs : $a,b$ valent 0 ou 2, et c vaut 0. Les quatre choix donnent $\res{1,4,9,36}$. Cette analyse épuise tous les exposants possibles, elle prouve donc la complétude de la liste.
\end{blocEdition}
\editionPage{defi-cor-4}{D04 · Corrigé}{L’identité de Sophie Germain}{\hyperref[defi-ex-4]{Énoncé p.\,\pageref*{defi-ex-4}}}
\begin{blocEdition}{defi-cor-4}
$(a^2+2b^2)^2-(2ab)^2=a^4+4a^2b^2+4b^4-4a^2b^2=a^4+4b^4$. Avec $a=n$ et $b=1$ : $n^4+4=(n^2+2)^2-(2n)^2=\res{(n^2-2n+2)(n^2+2n+2)}$. Or $n^2-2n+2=(n-1)^2+1\ge2$ pour $n\ge2$, et l’autre facteur est plus grand encore. Le nombre est le produit de deux entiers supérieurs ou égaux à 2 : $\res{\text{il n’est pas premier}}$. Pour $n=1$, on obtient 5, qui est premier. Cette identité porte le nom de Sophie Germain, mathématicienne française du XIX\textsuperscript{e} siècle.
\end{blocEdition}
\editionPage{defi-cor-5}{D05 · Corrigé}{Deux solutions et un contexte}{\hyperref[defi-ex-5]{Énoncé p.\,\pageref*{defi-ex-5}}}
\begin{blocEdition}{defi-cor-5}
La différence carré moins rectangle vaut $(x+1)^2-x(x+3)=1-x$. Une différence de 1 en grandeur signifie $1-x=1$ ou $1-x=-1$. On obtient x = 0 ou x = 2. La valeur 0 ne forme pas un rectangle de côté positif. Seule $\res{x=2}$ convient : rectangle d’aire 10 cm$^2$, carré d’aire 9 cm$^2$. Le contexte élimine une solution algébrique.
\end{blocEdition}
\editionPage{defi-cor-6}{D06 · Corrigé}{Promotions dans quel ordre ?}{\hyperref[defi-ex-6]{Énoncé p.\,\pageref*{defi-ex-6}}}
\begin{blocEdition}{defi-cor-6}
Réduction puis bon : $A=0{,}8P-10$. Bon puis réduction : $B=0{,}8(P-10)=0{,}8P-8$. Alors $B-A=\res{2\text{ €}}$. Le premier ordre est toujours plus avantageux de 2 €, pour les prix autorisés. Le bon de 10 € réduit avant pourcentage ne diminue finalement le prix que de 8 €. Les deux opérations ne commutent pas.
\end{blocEdition}
\editionPage{defi-cor-7}{D07 · Corrigé}{Une rencontre de débits}{\hyperref[defi-ex-7]{Énoncé p.\,\pageref*{defi-ex-7}}}
\begin{blocEdition}{defi-cor-7}
Pendant 12 minutes, le débit net vaut $18-3=15$ L/min. Le volume atteint est $15\times12=180$ L. La fuite vide ensuite 180 L à 3 L/min : il faut 60 minutes supplémentaires. Le temps total est $\res{72\text{ min}}$. On suppose la cuve assez grande et les débits constants ; aucun débordement n’est prévu.
\end{blocEdition}
\editionPage{defi-cor-8}{D08 · Corrigé}{Une fonction cachée derrière des essais}{\hyperref[defi-ex-8]{Énoncé p.\,\pageref*{defi-ex-8}}}
\begin{blocEdition}{defi-cor-8}
La fonction $g(x)=x^2+2+(x+1)x(x-1)$ possède bien les trois images indiquées, car le produit ajouté s’annule en -1, 0 et 1. Mais $f_0(2)=6$ pour la formule proposée et $g(2)=6+3\times2\times1=\res{12}$. Le tableau ne détermine donc pas une fonction unique. Une hypothèse supplémentaire sur sa famille serait nécessaire.
\end{blocEdition}
\editionPage{defi-cor-9}{D09 · Corrigé}{Trois tarifs, un choix argumenté}{\hyperref[defi-ex-9]{Énoncé p.\,\pageref*{defi-ex-9}}}
\begin{blocEdition}{defi-cor-9}
A = B pour n = 10 ; B = C pour $n=40/3\approx13{,}33$ ; A = C pour $n=35/3\approx11{,}67$. Pour 1 à 9 séances, A est la moins chère. Pour 10, A et B sont à égalité à 80 € et C est plus chère. Pour 11 à 13, B gagne. Pour 14 à 20, C gagne. Il faut comparer les trois coûts : l’intersection A-C ne signifie pas que ces deux offres soient les meilleures à cet endroit.
\end{blocEdition}
\editionPage{defi-cor-10}{D10 · Corrigé}{Une série à reconstruire}{\hyperref[defi-ex-10]{Énoncé p.\,\pageref*{defi-ex-10}}}
\begin{blocEdition}{defi-cor-10}
On impose $v_1=2$, $v_2=5$, $v_6=12$, $v_8=20$ et $v_4+v_5=16$. Pour une moyenne de 9, la somme vaut 72, donc $v_3+v_7=72-55=17$. La série $\res{2\,;5\,;5\,;7\,;9\,;12\,;12\,;20}$ convient : $Q_1=5$, médiane $(7+9)/2=8$, $Q_3=12$. Pour la seconde question, $v_3\le v_4\le8$, car $v_4\le v_5$ et $v_4+v_5=16$, et $v_7\le20$ : la somme vaut au plus $2+5+8+16+12+20+20=83$, donc la moyenne au plus $10{,}375$. Une moyenne de 15 est $\res{\text{impossible}}$.
\end{blocEdition}
\editionPage{defi-cor-11}{D11 · Corrigé}{Un jeu de somme équitable ?}{\hyperref[defi-ex-11]{Énoncé p.\,\pageref*{defi-ex-11}}}
\begin{blocEdition}{defi-cor-11}
Une somme 6 possède 5 couples, 8 en possède aussi 5, et 7 en possède 6. A gagne sur 10 couples, B sur 6. Parmi les 16 couples décisifs équiprobables, les parts sont $10/16=\res{5/8}$ et $6/16=\res{3/8}$. Le jeu avantage A. Recommencer les lancers indécis ne rend pas les trois sommes équiprobables.
\end{blocEdition}
\editionPage{defi-cor-12}{D12 · Corrigé}{Une corde dans un cercle}{\hyperref[defi-ex-12]{Énoncé p.\,\pageref*{defi-ex-12}}}
\begin{blocEdition}{defi-cor-12}
Le triangle OHA est rectangle en H. $HA^2=OA^2-OH^2=169-25=144$, donc HA = 12 cm. H étant le milieu de AB, $\res{AB=24\text{ cm}}$. L’aire vaut $24\times5/2=\res{60\text{ cm}^2}$. Le rayon n’est pas la hauteur relative à AB : c’est OH qui lui est perpendiculaire.\label{fig-defi-cor-12}\figureGeo[corrige]{defi-corde}

\end{blocEdition}
\editionPage{defi-cor-13}{D13 · Corrigé}{Les diagonales d’un trapèze}{\hyperref[defi-ex-13]{Énoncé p.\,\pageref*{defi-ex-13}}}
\begin{blocEdition}{defi-cor-13}
Les points A, O, C sont alignés, ainsi que B, O, D, et $(AB)\parallel(CD)$ : c’est un papillon de sommet O. Le théorème de Thalès donne $OA/OC=OB/OD=AB/CD=\res{6/10=3/5}$. Si AC = 12 cm, on partage 12 cm selon le ratio $3:5$ : $\res{OA=4{,}5\text{ cm}}$ et $\res{OC=7{,}5\text{ cm}}$. Contrôle : $4{,}5/7{,}5=0{,}6$.\label{fig-defi-cor-13}\figureGeo[corrige]{defi-trapeze}

\end{blocEdition}
\editionPage{defi-cor-14}{D14 · Corrigé}{Deux angles de visée}{\hyperref[defi-ex-14]{Énoncé p.\,\pageref*{defi-ex-14}}}
\begin{blocEdition}{defi-cor-14}
Notons x la distance entre B et le pied. L’angle de 45 degrés impose une hauteur h = x. Depuis A, $h/(x+10)=\tan30^\circ$. Donc $x(1-\tan30^\circ)=10\tan30^\circ$, puis $x=10\tan30^\circ/(1-\tan30^\circ)\approx13{,}6603$. La hauteur est $\res{13{,}7\text{ m}}$. Les distances et la hauteur obtenues sont positives.\label{fig-defi-cor-14}\figureGeo[corrige]{defi-visees}

\end{blocEdition}
\editionPage{defi-cor-15}{D15 · Corrigé}{Une somme de deux vecteurs}{\hyperref[defi-ex-15]{Énoncé p.\,\pageref*{defi-ex-15}}}
\begin{blocEdition}{defi-cor-15}
Par Chasles, $\vect{AE}=\vect{AB}+\vect{BE}=\vect{AB}+\vect{AC}=\vect{AD}$, donc $\res{E=D}$. Comme $\vect{BD}=\vect{AC}$, ABDC est $\res{\text{un parallélogramme}}$. Ses diagonales [AD] et [BC] ont le même milieu I, donc $\vect{AD}=2\vect{AI}$ : $\res{\vect{AB}+\vect{AC}=2\vect{AI}}$. Cette égalité est le point de départ de nombreuses démonstrations au lycée.\label{fig-defi-cor-15}\figureGeo[corrige]{defi-vecteurs}

\end{blocEdition}
\editionPage{defi-cor-16}{D16 · Corrigé}{Une section d’aire imposée}{\hyperref[defi-ex-16]{Énoncé p.\,\pageref*{defi-ex-16}}}
\begin{blocEdition}{defi-cor-16}
Le rayon du disque vaut 8 cm car $\pi r^2=64\pi$ et r est positif. Pythagore donne $OH^2=100-64=36$, donc $\res{OH=6\text{ cm}}$. Pour une direction de coupe fixée, deux plans parallèles conviennent, à 6 cm de part et d’autre du centre. Sans direction imposée, il existe de nombreuses orientations de plans.\label{fig-defi-cor-16}\figureGeo[corrige]{defi-section}

\end{blocEdition}
\editionPage{defi-cor-17}{D17 · Corrigé}{Un script de recherche et sa preuve}{\hyperref[defi-ex-17]{Énoncé p.\,\pageref*{defi-ex-17}}}
\begin{blocEdition}{defi-cor-17}
Après k tours, s est la somme de 1 à k et n vaut k+1. Les sommes passent par 91 pour k = 13 puis 105 pour k = 14. Le test devient faux après le quatorzième tour ; la sortie contient $\res{n=15}$ et s = 105. Une écriture par blocs utilise « répéter jusqu’à s $\ge$ 100 ». La somme augmente d’au moins 1 à chaque tour, ce qui garantit l’arrêt.
\end{blocEdition}
\editionPage{defi-cor-18}{D18 · Corrigé}{Le meilleur rectangle entier}{\hyperref[defi-ex-18]{Énoncé p.\,\pageref*{defi-ex-18}}}
\begin{blocEdition}{defi-cor-18}
L’aire vaut $x(23-x)$ avec x entier entre 1 et 22. On calcule $132-x(23-x)=(x-11)(x-12)$. Pour un entier x, les facteurs sont tous deux négatifs ou nuls si $x\le11$, tous deux positifs ou nuls si $x\ge12$. Leur produit est donc positif ou nul : l’aire ne dépasse pas 132. L’égalité est atteinte pour x = 11 ou 12, soit le même rectangle de $\res{11\text{ cm sur }12\text{ cm}}$. L’argument utilise le caractère entier : aucun entier n’est strictement entre 11 et 12.
\end{blocEdition}
